\section*{Bloque C. Casos con implicancias tarifarias y regulatorias del sector saneamiento}

\noindent\textit{Preguntas 66 a 90. Aplican los conceptos de competencia perfecta, equilibrio de mercado, excedentes, incidencia de impuestos y controles de precios a mercados de agua, servicios de saneamiento y políticas tarifarias.}

\q{En un mercado competitivo de camiones cisterna, la demanda es $Q_d=80-2P$ y la oferta es $Q_s=2P-20$. El precio de equilibrio es:}
\begin{opts}
    \item 20
    \item 25
    \item 30
    \item 35
\end{opts}

\q{En el mercado anterior, el excedente del consumidor es:}
\begin{opts}
    \item 112.5
    \item 225
    \item 450
    \item 900
\end{opts}

\q{En el mismo mercado, el excedente del productor es:}
\begin{opts}
    \item 112.5
    \item 450
    \item 900
    \item 225
\end{opts}

\q{En el mismo mercado, el excedente total es:}
\begin{opts}
    \item 225
    \item 900
    \item 450
    \item 675
\end{opts}

\q{Si el regulador fija un precio máximo de 20 en ese mercado, la escasez es:}
\begin{opts}
    \item 10
    \item 20
    \item 30
    \item 40
\end{opts}

\q{Con el precio máximo de 20, la pérdida de bienestar es:}
\begin{opts}
    \item 50
    \item 75
    \item 200
    \item 100
\end{opts}

\q{Si se aplica un impuesto específico de 4 por unidad, el nuevo precio de equilibrio es:}
\begin{opts}
    \item 25
    \item 27
    \item 29
    \item 31
\end{opts}

\q{Con el impuesto de 4, la recaudación total es:}
\begin{opts}
    \item 52
    \item 80
    \item 100
    \item 104
\end{opts}

\q{Con el impuesto de 4, la pérdida de bienestar es:}
\begin{opts}
    \item 8
    \item 16
    \item 4
    \item 12
\end{opts}

\q{En el impuesto anterior, la carga que soportan los consumidores por unidad es:}
\begin{opts}
    \item 1
    \item 3
    \item 4
    \item 2
\end{opts}

\q{En el mismo impuesto, la carga que soportan los productores por unidad es:}
\begin{opts}
    \item 1
    \item 2
    \item 3
    \item 4
\end{opts}

\q{Si se aplica un subsidio de 4 por unidad, el nuevo precio de equilibrio es:}
\begin{opts}
    \item 21
    \item 25
    \item 23
    \item 27
\end{opts}

\q{Con el subsidio de 4, el costo total del subsidio es:}
\begin{opts}
    \item 104
    \item 120
    \item 136
    \item 160
\end{opts}

\q{Con el subsidio de 4, la pérdida de bienestar es:}
\begin{opts}
    \item 8
    \item 16
    \item 4
    \item 12
\end{opts}

\q{Si en una industria competitiva el precio es mayor que el costo medio mínimo, en el largo plazo:}
\begin{opts}
    \item salen empresas de la industria y el precio tiende a subir.
    \item entran empresas, sube la oferta y el precio tiende a bajar.
    \item el equilibrio de mercado permanece completamente inalterado.
    \item el precio sigue subiendo de forma indefinida en el tiempo.
\end{opts}

\q{En una industria de costo constante, la oferta de largo plazo es:}
\begin{opts}
    \item horizontal.
    \item ascendente.
    \item descendente.
    \item vertical.
\end{opts}

\q{En un monopolio natural con costo medio decreciente, fijar el precio igual al costo marginal genera:}
\begin{opts}
    \item beneficios extraordinarios positivos.
    \item beneficio económico exactamente nulo.
    \item ausencia total de déficit financiero.
    \item pérdidas que exigen subsidio o cargo.
\end{opts}

\q{La teoría de mercados contestables sostiene que, incluso con pocas empresas, el precio puede aproximarse al competitivo si:}
\begin{opts}
    \item existen barreras legales infranqueables a la entrada.
    \item la amenaza de entrada de nuevos rivales es creíble.
    \item la empresa incumbente actúa como monopolio puro cerrado.
    \item las empresas establecidas fijan precios coordinadamente.
\end{opts}

\q{Si la tarifa de agua disminuye y todo lo demás permanece constante, el excedente del consumidor:}
\begin{opts}
    \item disminuye.
    \item aumenta.
    \item no cambia.
    \item se vuelve cero.
\end{opts}

\q{Si la tarifa de agua disminuye y todo lo demás permanece constante, el excedente del productor:}
\begin{opts}
    \item aumenta proporcionalmente con el volumen de agua.
    \item permanece exactamente igual en el corto plazo.
    \item disminuye y puede comprometer la sostenibilidad.
    \item se vuelve nulo aunque el precio siga sobre $CVMe$.
\end{opts}

\q{En un mercado de derechos de agua, la asignación competitiva tiende a:}
\begin{opts}
    \item repartir el agua por igual entre todos los usuarios.
    \item asignar el agua a los usos de mayor valor económico.
    \item asignar el agua a los usos de menor valor.
    \item distribuir el agua de forma aleatoria.
\end{opts}

\q{Un impuesto pigouviano sobre el consumo de agua que genera un costo externo debe fijarse igual a:}
\begin{opts}
    \item cero.
    \item el costo externo marginal.
    \item el costo medio total.
    \item el costo total del servicio.
\end{opts}

\q{Un precio máximo por debajo del equilibrio en un mercado de agua genera:}
\begin{opts}
    \item un excedente persistente de oferta.
    \item ningún efecto sobre las transacciones.
    \item escasez y racionamiento sin precios.
    \item un aumento automático de la oferta.
\end{opts}

\q{Si el objetivo es proteger a los hogares vulnerables sin generar escasez, la opción más eficiente es:}
\begin{opts}
    \item fijar un precio máximo general por debajo del equilibrio.
    \item establecer un precio mínimo obligatorio superior al equilibrio.
    \item imponer cuotas físicas uniformes de consumo por vivienda.
    \item aplicar un subsidio directo focalizado sin distorsionar precios.
\end{opts}

\q{El equilibrio competitivo se caracteriza por ser eficiente en sentido asignativo, pero:}
\begin{opts}
    \item garantiza siempre una distribución totalmente equitativa.
    \item resulta ineficiente incluso bajo información simétrica.
    \item prescinde de cualquier política pública redistributiva.
    \item puede requerir transferencias para lograr metas de equidad.
\end{opts}
